\scanpage{96}
plus une question pour, des courbes relatives
\uline{globales}, mais pour des ``éléments de
courbe'' relatifs --- savoir, pour le
complété formel d'une courbe relative
$X/B$ le long d'une section $\unsure{g_i|g}$,
avec le diviseur de cette section \unsure{distingu|disting} % word meets the right paper edge; no ending supplied
dans la structure.

Je vais plutôt revenir à la situation
où on part de une section de première
espèce d'un $\mathcal E_{U_K}$ sur $\mathcal E_K$, associée
donc à un pt \unsure{rationnel|rationnelle} $f_K$ de $U_K$ sur $K$.
On peut donc dire (indépendamment de
la donnée du modèle $S$ de $K$ sur
lequel $U_K$ \unsure{pouvait|pourrait} déjà se définir
en \unsure{tant|\deleted{tant}} que schéma relatif élémen% source line-break hyphen
taire) que la dite section de
$\mathcal E_{U_K}$ sur $\mathcal E_K$ \inserted{$(29)\quad\unsure{(f_K)_*|\mathcal E f_K|f_K}=\psi:\mathcal E_K\longrightarrow\mathcal E_{U_K}$} définit, pour \unsure{tt|tte} section
de $\mathcal E_K$ sur $\mathcal E_{\mathbf Q}=\Gamma_{\mathbf Q}$ ou sur un sous-%
groupe ouvert \inserted{$\Gamma'$} de celui-ci, une section
\margin{95}
